\chapter{State-free normal conductors at the large-$L$ corner}
\label{component:p3be5a8005801}
\noindent\textbf{Component }\path{p3be5a8005801}\par
{\small\noindent Source: \path{BASE/proofs/c1_unbounded/c141a_closure_and_c141b_partial.tex}\par}
{\small\noindent SHA256: \path{3be5a8005801fb7deeadb216e818574e376e5eaf56659c8337df32734e75bc6b}\par}
{\small\noindent Inventory use: SF1-SF9 source bounds, exact radial and normal allocation, surviving h-residue<1/2 C14.1b subdomain | SF1--SF18 and both theorem proofs | C14\_\allowbreak{}ENTIRE\_\allowbreak{}ORIGINAL\_\allowbreak{}PARENT\_\allowbreak{}GRAPH; precise subordinate uses retained below\par}
\medskip
\section{State-free normal conductors at the large-$L$ corner}
\label{p3be5a8005801:sec:c141-state-free}

\paragraph{Claims and retained input.}
Keep the original physical stationary C1 source, its positive $d$,
strict gaps, nonnegative event coefficients, original infinitesimal
total scale, and the complement of the previously admitted domains.
Normalize $d=1$. Source quantities belong to a real closed extension
$K_c$ and testing angles and states to an independent extension $K_t$.
The conclusion is a finite sum of squares in
$K_c\otimes_{\mathbb R}K_t$, not pointwise positivity after identifying
the fields. In this section $x,z=1/L,\sigma=1/\eta$ are all positive
infinitesimals. All original source openings are retained, including
the allowed literal faces $D=0$ and $\beta=0$.

The first theorem closes the entire original C14.1a residual. The
second closes a proper part of C14.1b. The precise remaining C14.1b
complement is stated at the end. No global theorem is claimed.

\begin{theorem}[C14.1a]
If $H=h/\sigma$ is bounded, the complete original parent is tensor
SOS for every independent projective test and every real or complex
state. In particular the previously unresolved C14.1a domain,
$M_\beta=\beta\theta/\sigma^2$ unbounded, is admitted completely.
\end{theorem}

\begin{theorem}[A further C14.1b subdomain]
If $H$ is unbounded and the real residue of $h$ is strictly less
than $1/2$, the same full-parent conclusion holds. This includes
every infinitesimal $h$ and every positive unit $h$ with residue
less than $1/2$, but does not include residue $1/2$ or greater.
\end{theorem}

The proof uses the exact complete P4--P8/SP6 parent in the inherited
\path{parent_reference.tex}, the BE1--BE9 source identities, and the
already admitted remote and literal-angle exits. The new argument
settles the remaining $u,w=1/v$ infinitesimal corner directly; it does
not reclassify any source outside the original C14.1 domains.

\subsection{Source bounds with $z$ infinitesimal}

Write $q=\sqrt z$, $k=\sqrt x$, $b=Lr$, $Q=1-b$,
$\theta=1-h$, $e=xh/\sigma$, and retain the exact source letters
$J,M,P,B,\mathcal H,p,l,\kappa$ from P1--P2. Since $r=bz$ and
$0<b<1$, $r$ is infinitesimal and $1-r$ is a unit. BE1 applies
because it requires only that $L$ have a positive real lower bound;
it gives $e=O(x+\sigma)$ and $E=1-e\sim1$. Consequently
\[
 P,J,M,B,\mathcal H\sim1,\qquad
 p,l\sim\tfrac12,\qquad C=O(z),\qquad\kappa\sim\tfrac12.
 \tag{SF1}
\]
The possible $b^{-1},Q^{-1}$ poles occur only in the two original
positive source weights $a=Az=p/b$ and $\ell=\alpha z=l/Q$.
Both have positive real lower bounds. They need not be bounded above.
No cap of an individual unbounded matrix entry is used below.

Let $d_0,b_0$ be the exact source sphere from P2, and put
\[
 d=4D/\sigma^2,\qquad
 \lambda=2H+\lambda_0,\qquad
 \rho=\lambda x/\sigma,\qquad
 K_A=4\kappa H-\lambda.
 \tag{SF2}
\]
The identities SP9--SP10 show that $\nu,\lambda_0,\chi,K_A$ are
bounded even when $H$ is unbounded. This is a polynomial
divisibility argument, not multiplication of an uncontrolled error
by $H$. More precisely their source residues give
\[
 (\nu)_0=1,\qquad (\lambda_0)_0=-2(1-h_0),\qquad
 (K_A)_0=2-4h_0.
 \tag{SF3}
\]
For example the SP10 polynomial $F_\sigma$ has residue $-1$ at
$x=z=\sigma=0$, whereas $eF_x$ has zero residue. Thus
$\alpha_\kappa=2H(2\kappa-1)$ has residue $-2h_0$, which proves
the last equality without assuming a rate for $2\kappa-1$.

The source-sphere calculations HG3--HG5 are algebraic for bounded
$z$ and do not require it to be a unit. They give
\[
 \rho\text{ is bounded},\qquad
 (\rho)_0=2(b_0^2)_0,\qquad (d+\rho)_0=2.
 \tag{SF4}
\]
One way to check the delicate assertion is to use both exact identities
\[
 \rho=\frac D{\sigma^2}(8hx-4\theta J)
       +\frac{2xh^2}{\sigma^2}
       +\frac{2\theta\mathcal H}{P}-4bpxz,
 \qquad
 \frac{xh^2}{\sigma^2}=\frac{eH_0}{B}
                         +\frac{h\mathcal H}{B}b_0^2.
 \tag{SF5}
\]
These include $D=0$ and $\beta=0$. Throughout the positive-$\lambda$
case $\lambda x=\rho\sigma$ is infinitesimal.

When $H$ is bounded, $h$ is infinitesimal, $K_A\sim2$, and
$\lambda$ is bounded. If $\lambda\le0$, its exact defining formula
implies
\[
 8Dh+4\beta\theta^2+2h^2\le4bp\sigma^2z.
\]
Since $\theta\sim1$, this forces bounded $M_\beta$; the inherited
JB bounded-$M_\beta$ all-tests theorem applies. On the actual
C14.1a remainder $M_\beta$ is unbounded, so $\lambda>0$.

If $d_0$ is infinitesimal, including literal $D=0$, SF4 makes
$\rho$ a positive bounded unit. The positive sphere identity
\[
 H=H_0/B+(\sigma/x)(\mathcal H/B)b_0^2
 \tag{SF6}
\]
then gives the source bounds needed for the two theorems:
\[
 \begin{array}{ll}
 H\text{ bounded}:&\sigma/x\text{ bounded};\\
 H\text{ unbounded}:&x/\sigma\text{ infinitesimal},\quad
                    xh/\sigma^2\text{ a positive unit},\quad
                    \sigma/k\asymp\sqrt h\text{ bounded}.
 \end{array}
 \tag{SF7}
\]
For the second line use $\lambda/H\sim2$ and $\rho\sim2$.
None of these is a comparison between a source rate and a testing rate.

\subsection{The full source-opening-safe radial block}

The roots are the original ones:
\[
 N_x=X+uhZ,\quad N_y=(Y-bwZ)/q,\quad
 \Gamma=\sigma Z-uwW,\quad R=wZ/q-uqW,\quad
 A_\theta=q(u^2-x)W.
\]
Put $J_C=N_y+bR$ and $J_\ell=N_y-QR$. The exact CT7--CT7a
radial extraction, with every $a,\ell$ excess restored, has leading
pure-source form
\[
 [4a+2z(2C+1)]J_C^2+4\ell J_\ell^2.
 \tag{SF8}
\]
Retain the nonconstant positive $u^2,w^2$ packets and all positive
source-weight excesses. Clipping the two coefficients in SF8 at
positive real lower bounds gives a strict bounded Gram on $(N_y,R)$,
since
\[
 R=J_C-J_\ell,\qquad N_y=QJ_C+bJ_\ell,
 \qquad b+Q=1.
\]
This is uniform through $b\to0$ and $Q\to0$. The remaining three
radial coefficients are precisely CT7a. They are infinitesimal
because $z$ is infinitesimal, $C=O(z)$, and all their other source
letters are bounded. In particular the $z$-unit argument's possible
$C$ pole does not occur here.

The $N_x^2$ diagonal has residue $2$. The complete remaining
non-Gamma cross terms of SP6 are
\[
 \begin{split}
 &-4uqw(2m+2D+1)N_xN_y+uqwR_xN_xR\\
 &\hspace{4mm}+qwB_xN_xA_\theta+uA_yN_yA_\theta
                         +\chi uA_\theta R.
 \end{split}
 \tag{SF9}
\]
Every coefficient is a bounded infinitesimal on $u,w$ infinitesimal.
Here $m,D,C,\beta\theta,p,l$ are bounded, $B_x=8\kappa e$ is
infinitesimal, and $A_y$ is bounded. The entire $N_yR$ term was
already included in SF8 and CT7a. No first normal was suppressed.

\subsection{A scalar circuit and two normal conductors with no state quotient}

The original-core sign exit and exact testing sign gauges leave
$u,w,W,Z>0$. Literal $u=0,w=0,W=0,Z=0$ and the opposite-sign
cases have their inherited direct core SOS. Suppose $\lambda>0$.
Define separated nonnegative monomials
\[
 a_c=\sqrt{\lambda\sigma}\,uZ,\qquad
 b_c=\sqrt\rho\,uwW,\qquad
 c_c=\sqrt\lambda\,xqW,\qquad
 g_c=(a_cb_cc_c)^{1/3},\qquad
 s_c^3=\frac{wZ}{zuW}.
 \tag{SF10}
\]
Each radical is taken separately in its source and testing factors.
The exact full angular packet in SP6 is
\[
 K_AA_\theta^2+\chi uA_\theta R+\mathscr P,
 \qquad
 \mathscr P=\lambda\sigma P_\beta^2
       +\rho(u^2-x)(w^2-\sigma z)W^2.
\]
The same algebraic HP3 identity is valid for infinitesimal $z$:
\[
 \mathscr P+\lambda xR^2
 =a_c^2+b_c^2+c_c^2-3g_c^2
   +\lambda xzu^2W^2s_c^2(s_c-1)^2(s_c^2+2s_c+3).
 \tag{SF11}
\]
No bound on $s_c$ or any state ratio is asserted or required. The last
term is a legal SOS after writing $s_c^2+2s_c+3=(s_c+1)^2+2$.
Furthermore
\[
 a_c^2+b_c^2+c_c^2-3g_c^2
 =\tfrac12\{(a_c-b_c)^2+(b_c-c_c)^2+(c_c-a_c)^2\}
      +(a_cb_c+b_cc_c+c_ca_c-3g_c^2).
 \tag{SF12}
\]
The last parenthesis is the elementary cubic AM--GM SOS applied to
$(a_cb_c)^{1/3},(b_cc_c)^{1/3},(c_ca_c)^{1/3}$. Expand its
positive monomial prefactor before taking square roots. Thus
\[
 D_1=(a_c-b_c)/\sqrt2,\qquad
 D_2=(b_c-c_c)/\sqrt2
\]
are two genuine unit-weight pivots; all other SF11--SF12 squares
are retained separately.

The new root is
\[
 T=kq(u-k)W.
\]
Its legal quadratic domination is exact:
\[
 A_\theta^2-T^2=q^2(u-k)^2W^2(u^2+2uk).
 \tag{SF13}
\]
Both summands on the right are squares, using the separated root
$\sqrt u\sqrt k$. If $K_A$ has positive residue, allocate its
diagonal exactly once as
\[
 K_AA_\theta^2=\tfrac12K_AA_\theta^2+\tfrac12K_AT^2
                        +\tfrac12K_A(A_\theta^2-T^2).
 \tag{SF14}
\]
The last term is retained SOS and the first two supply strict
positive source pivots for $A_\theta,T$.

Here are the decisive two identities, valid for arbitrary states
without a state division:
\[
 \boxed{\begin{aligned}
 u\Gamma={}&\frac{\sqrt2 k}{\sqrt\rho}D_1
       -\frac{\sqrt2(u-k)}{\sqrt\rho}D_2-\sqrt\sigma T,\\
 qw\Gamma={}&\sigma zR
       -\frac{\sqrt2(qw+z\sqrt\sigma)}{\sqrt\rho}D_2
       +\left(\frac{\sigma z}{k}
                  +\frac{\sqrt\sigma qw}{k}\right)T.
 \end{aligned}}
 \tag{SF15}
\]
All divisions are pure source. In particular neither $q$ nor a
testing coordinate is inverted in a conductor coefficient. These
replace the $\lambda(Z/W)/q$ and $w/q$ coefficients that obstruct a
literal transport of the R19 $z$-unit proof.

\subsection{Completion when $H$ is bounded}

Here $K_A\sim2$ and $\lambda$ is bounded. If $d_0$ has positive
residue, the original $d\Gamma^2$ diagonal gives a strict Gamma
pivot. Its two complete normals are
\[
 -4\nu N_x(u\Gamma)
       +R_\gamma(N_y+bR)(qw\Gamma).
 \tag{SF16}
\]
Their coefficients are infinitesimal on $(N_x,N_y,R,\Gamma)$.
Keep SF11 entirely as SOS; its sole cost $-\lambda xR^2$ is an
infinitesimal bounded radial correction. The frame
$(N_x,N_y,R,\Gamma,A_\theta)$ has strict positive residue.

If $d_0$ is infinitesimal, $\rho$ is a unit by SF4 and
$\sigma/x$ is bounded by SF7. All coefficients in both lines of
SF15 are bounded infinitesimals. Indeed $\sqrt\sigma/k$ is
bounded, $\sigma/k=k(\sigma/x)$ is infinitesimal, and $q,w,z$
are bounded infinitesimals. Use SF14 and the two unit circuit
pivots. After substituting SF15 into SF16, the complete frame
\[
 (N_x,N_y,R,A_\theta,T,D_1,D_2)
 \tag{SF17}
\]
has strictly positive residue. Its radial block is SF8, its angular
and $T$ pivots come from SF14, and the circuit pivots from SF12.
Every correction is a bounded infinitesimal. The diagonal
$d\Gamma^2$ is retained separately, including at $d=0$.

In both cases retain exactly once $4\beta(1+w^2)N_x^2$,
$2(2D+1)(uX_e)^2$, $4bpzu^2\Gamma^2$, the unused circuit rows,
and all positive radial and coefficient-cap excesses. Equations
SP6, SF8--SF16 list every original scalar, normal and hard coupling.
The separated bounded-Gram lemma proves the entire parent.

The $d_0$ split exhausts the source sphere. The bounded-$M_\beta$
and bounded nonzero-unit Alice-angle exits are the inherited JB
theorems; their CE remote join includes all other projective tests.
Together with the literal-axis and sign exits, this closes the
complete original C14.1a residual and the full bounded-$H$ source
domain. Apply the real argument independently to real and imaginary
state parts and restore only positive pure-source and pure-testing
normalizations.

\subsection{A source factor restores the second normal when $H$ is unbounded}

Assume now $H$ unbounded and $h_0<1/2$. SF3 gives a positive-unit
$K_A$, so the same angular allocation is available. Also
$\lambda>0$, $\lambda/H\sim2$, $\rho$ bounded and
$\lambda x=\rho\sigma$ infinitesimal. The positive-$d_0$ case is
closed exactly as before: $\lambda$ occurs only inside the retained
scalar circuit, not as an unbounded coefficient of the remaining Gram.

For infinitesimal $d_0$, $\rho$ is a unit, but $\sqrt\sigma/k$
can be unbounded. One must use the actual second normal in SF16,
not claim boundedness of its factor $qw\Gamma$ by itself. Exactly
\[
 R_\gamma=R_0-8pu^2,\qquad
 R_0=-4\{2(C+D)(r-1)+2p+r-1\}.
 \tag{SF18}
\]
The required source bound is $R_0/\sigma$ bounded. Here is a
reproducible polynomial proof that retains every opening. In the
P1 formulas temporarily leave $e$ independent, and put
$N=JP R_0$. The cancellation of $1-r$ in $C(r-1)$ shows that $N$
is a polynomial in $x,\sigma,z,b,h,e$. Moreover
$N|_{x=\sigma=e=0}=0$. Define the polynomial quotients
\[
 F_x=\frac{N-N|_{x=0}}x,\quad
 F_\sigma=\frac{N|_{x=0}-N|_{x=\sigma=0}}\sigma,\quad
 F_e=\frac{N|_{x=\sigma=0}}e.
\]
All their source arguments are bounded and the denominators $JP$
are units. Restoring $e=xh/\sigma$ gives the exact expression
\[
 \frac{R_0}{\sigma}
 =\frac{(x/\sigma)F_x+F_\sigma+(xh/\sigma^2)F_e}{JP}.
 \tag{SF19}
\]
The two ratios on the right are bounded by SF7; this proves the
claim. The checker prints the three actual polynomials in its receipt
so that this is not an unverifiable divisibility assertion.

Use the first conductor in SF15 unchanged. For the entire other
normal use the exact identity
\[
 qwR_\gamma\Gamma=R_0(qw\Gamma)-8p qwu(u\Gamma).
 \tag{SF20}
\]
Substitute both SF15 expressions into SF20. Every resulting
coefficient is infinitesimal. The only apparently dangerous ones are
\[
 R_0\frac{\sigma z}{k},\qquad
 R_0\frac{\sqrt\sigma qw}{k}.
\]
Write $R_0=\sigma C_0$ with bounded pure-source $C_0$. SF7 makes
$\sigma/k$ bounded, so these are respectively
$C_0\sigma z(\sigma/k)$ and
$C_0\sqrt\sigma qw(\sigma/k)$, both infinitesimal. The other
term in SF20 has the explicit infinitesimal factor $qwu$.

Thus SF17 again has strict positive residue, now using SF20 for the
complete second Gamma normal. All retained squares and every other
SP6 term are accounted for as in the bounded-$H$ proof. The inherited
R17 bounded-state and R18 unit-Alice admissions, together with the
unchanged remote joins, complete all independent projective tests on
this stated source subdomain. The argument never inverts a sphere
coordinate and remains literal at $D=0$.

\subsection{Exact remaining complement and verification boundary}

C14.1a has no remaining mathematical domain after this proof is
combined with the recorded inherited testing exits. C14.1b remains
open on its actual inherited complement intersected with
\[
 H\text{ unbounded},\qquad \operatorname{res}(h)\ge\tfrac12.
 \tag{SF21}
\]
The earlier bounded-$W/Z$, sign, literal-axis and nonzero-unit-$u$
admissions still remove their stated parts. Thus its real residual
still has $u,w>0$ infinitesimal, $WZ>0$ and $W/Z$ unbounded,
as well as every original source intersection.

The specific obstruction to extending this proof is the angular
diagonal: $K_A$ has residue $2-4h_0$, which is zero or negative on
SF21. Consequently SF14 supplies no strict positive angular or $T$
pivot there. A new allocation of this signed angular packet is
required. This failed frame is not a counterexample to tensor SOS
and does not establish that SF21 is nonempty.

The focused file \path{check_state_free_conductors.py} checks the
two state-free identities, the entire second normal substitution,
the exact once-paid angular budget, the complete scalar circuit,
its pair-difference pivots, and the source cancellation in SF19.
It reports fifteen exact assertions. The complete original SP6 and
CT7a identities remain inherited checked inputs. The analytical
argument above, not the assertion count, establishes the quantified
admissions. Review is same-system mathematical review, not external
peer review or formal verification.
